% TOP_PROBLEM: {"problemId": "problem.prove-the-penrose-inequality-or-present-a-counterexample", "canonicalTitle": "Prove the Penrose inequality or present a counterexample", "releaseRank": 37, "primaryDomain": "mathematical_physics", "primaryDomainLabel": "Mathematical physics", "displayStatus": "open_with_solved_subcases", "releaseStatus": "open_with_solved_subcases"}
% !TeX program = lualatex
% Definition and sources only; this file is not an LLM solution attempt.
\documentclass[11pt]{article}
\usepackage[margin=25mm]{geometry}
\usepackage{fontspec}
\setmainfont{Latin Modern Roman}
\usepackage{amsmath,amssymb,mathtools}
\usepackage{unicode-math}
\setmathfont{Latin Modern Math}
\usepackage{xurl,hyperref}
\hypersetup{colorlinks=true,urlcolor=blue}
\setcounter{secnumdepth}{0}
\setlength{\parindent}{0pt}
\setlength{\parskip}{5pt}
\setlength{\emergencystretch}{3em}
\begin{document}
\section*{\texorpdfstring{Prove the Penrose inequality or present a counterexample}{Prove the Penrose inequality or present a counterexample}}\label{37}
\subsection{Definitions and mathematical statement}
An initial data set consists of a Riemannian $3$-manifold $(M,g)$ and symmetric tensor $k$, its second fundamental form. In geometric units the constraint densities obey
\[
 16\pi\mu=R_g+(\operatorname{tr}_gk)^2-|k|_g^2,\qquad
 8\pi J=\operatorname{div}_g(k-(\operatorname{tr}_gk)g).
\]
The dominant energy condition is $\mu\ge|J|_g$. At an asymptotically Euclidean end, ADM energy is defined by the limiting metric-flux integral; the mass convention and decay hypotheses are those of the cited formulation. If $A$ is the infimum of areas enclosing the specified outermost apparent horizon, the desired bound is
\[
 m_{\rm ADM}\ge\sqrt{A/(16\pi)},
\]
with the source's Schwarzschild rigidity condition. An apparent horizon is characterized by vanishing appropriate outgoing null expansion. The full spacetime statement does not assume $k=0$.
\par\smallskip\textit{Formulation and concept references:} \href{https://arxiv.org/abs/2512.04137}{[S1]}.

\subsection{Short English statement}
Does the mass of an asymptotically flat spacetime initial state always dominate the area-based lower bound from its outermost horizon?
\subsection{Sources}
\begin{itemize}
\item[S1]D. Xu, 'The Spacetime Penrose Inequality: Conditional Results for Stable MOTS and General Trapped Surfaces,' arXiv:2512.04137v5 (2025).
\url{https://arxiv.org/abs/2512.04137}.
\textit{Formulation source.}
\item[S2]The spacetime Penrose inequality under a quasi final state hypothesis.
\url{https://arxiv.org/abs/2605.18730}.
\textit{Further reference; not a proof endorsement.}
\item[S3]Proof of the Spacetime Penrose Inequality With Suboptimal Constant in the Asymptotically Flat and Asymptotically Hyperboloidal Regimes.
\url{https://arxiv.org/abs/2504.10641}.
\textit{Further reference; not a proof endorsement.}
\end{itemize}
\end{document}
